\subsection{Haar measure of a product}

PROPOSITION 5. --- Let \((\mathbf{G}_{\iota})_{\iota \in \mathbf{I}}\) be a family of locally compact groups. For every \(\iota \in \mathbf{I}\) let \(\mu_{\iota}\) be a left (resp. right) Haar measure on \(\mathbf{G}_{\iota}\) . Assume that there exists a finite subset \(\mathbf{J}\) of \(\mathbf{I}\) such that, for every \(\iota \in \mathbf{I} - \mathbf{J}\) , \(\mathbf{G}_{\iota}\) is compact and \(\mu_{\iota}(\mathbf{G}_{\iota}) = 1\) . Then the product measure \(\bigotimes_{\iota \in \mathbf{I}} \mu_{\iota}\) is a left (resp. right) Haar measure on \(\mathbf{G} = \prod_{\iota \in \mathbf{I}} \mathbf{G}_{\iota}\) . If \(x = (x_{\iota}) \in \mathbf{G}\) then

\[
\Delta_ {\mathrm{G}} (x) = \prod_ {\iota \in \mathrm{I}} \Delta_ {\mathrm{G} _ {\iota}} (x _ {\iota}).
\]

For every finite subset \(J\) of \(I\) , \(\bigotimes_{\iota \in J} \mu_{\iota}\) is a left (resp. right) Haar measure on \(\prod_{\iota \in J} G_{\iota}\) , as follows immediately from the definitions. Therefore \(\bigotimes_{\iota \in I} \mu_{\iota}\) is a left (resp. right) Haar measure on \(G\) (Ch. III, §4, No.~6, Prop. 9). On the other hand, if the \(\mu_{\iota}\) are left Haar measures then

\[
\boldsymbol {\delta} (x) \left(\bigotimes_ {\iota \in I} \mu_ {\iota}\right) = \bigotimes_ {\iota \in I} \boldsymbol {\delta} \left(x _ {\iota}\right) \mu_ {\iota} = \bigotimes_ {\iota \in I} \left(\Delta_ {G _ {\iota}} \left(x _ {\iota}\right) \mu_ {\iota}\right) = \left(\prod_ {\iota \in I} \Delta_ {G _ {\iota}} \left(x _ {\iota}\right)\right) \bigotimes_ {\iota \in I} \mu_ {\iota},
\]

whence \(\Delta_{\mathbf{G}}(x) = \prod_{\iota \in \mathrm{I}}\Delta_{\mathbf{G}_{\iota}}(x_{\iota})\)

Examples. --- 1) Lebesgue measure on \(\mathbf{R}^n\) is a Haar measure of the additive group \(\mathbf{R}^n\) .

\begin{enumerate}
\def\labelenumi{\arabic{enumi})}
\setcounter{enumi}{1}
\tightlist
\item
  The mapping \((r, u) \mapsto ru\) is an isomorphism of \(\mathbf{R}_+^* \times \mathbf{U}\) onto \(\mathbf{C}^*\) (GT, VIII, §1, No.~3). If \(\mathbf{C}^*\) is identified with \(\mathbf{R}_+^* \times \mathbf{U}\) by means of this isomorphism, and if \(du\) denotes a Haar measure on \(\mathbf{U}\) , then \(r^{-1} dr du\) is a Haar measure on \(\mathbf{C}^*\) by Example 2 of No.~2. On the other hand, the bijection \(\theta \mapsto e^{2i\pi\theta}\) of {[}0, 1{[} onto \(\mathbf{U}\) transforms the Lebesgue measure \(d\theta\) on {[}0, 1{[} into a Haar measure on \(\mathbf{U}\) by Example 3 of No.~2. It follows that if \(f \in \mathcal{K}(\mathbf{C}^*)\) , the integral
\end{enumerate}

\[
\int_ {0} ^ {+ \infty} \int_ {0} ^ {1} f (r e ^ {2 i \pi \theta}) r ^ {- 1} d r d \theta
\]

defines a Haar measure on \(\mathbf{C}^*\) .

